% ==STYLE-BEGIN==
\documentclass[UTF8,11pt]{ctexart}
\usepackage{amsmath,amssymb,amsthm}
\usepackage[a4paper,margin=2.4cm,top=2.3cm,bottom=2.3cm]{geometry}
\usepackage{fancyhdr}
\usepackage{enumitem}
\usepackage{hyperref}
\hypersetup{hidelinks}

% ---- 版式 ----
\linespread{1.15}
\setlength{\parskip}{0.35em}
\setlength{\parindent}{1.6em}
\setlist{nosep,leftmargin=2em}

% ---- 页眉 ----
\pagestyle{fancy}
\fancyhf{}
\fancyhead[L]{\small\coursename}
\fancyhead[R]{\small\thepage}
\renewcommand{\headrulewidth}{0.4pt}

% ---- 定理环境：中文用正体，标题加粗，按小节共用一条流水号 ----
\newtheoremstyle{cnthm}{0.7em}{0.7em}{\normalfont}{}{\bfseries}{}{0.5em}{}
\theoremstyle{cnthm}
\newtheorem{thm}{定理}[section]
\newtheorem{lem}[thm]{引理}
\newtheorem{prop}[thm]{命题}
\newtheorem{cor}[thm]{推论}
\newtheorem{defn}[thm]{定义}
\newtheorem{rem}[thm]{注}
\newtheorem{exa}[thm]{例}
\newtheorem{conj}[thm]{猜想}
\renewcommand{\proofname}{\bfseries 证明}
\providecommand{\coursename}{}
% ==STYLE-END==
\renewcommand{\coursename}{数学分析\,·\,0924}
\title{数学分析 0924}
\date{}
\begin{document}
\maketitle

\section{Recall：有界函数}

\[
\exists M>0,\ \forall x\in X,\ |f(x)|<M
\]
称函数 $f(x)$ 有界。函数无界即
\[
\forall M>0,\ \exists x\in X,\ |f(x)|\geq M .
\]

对于无界函数，取 $M=1$，$\exists x_1\in X,\ |f(x_1)|\geq1$；
取 $M=2$，$\exists x_2\in X,\ |f(x_2)|\geq2$；
$\cdots$；取 $M=k$，$\exists x_k\in X,\ |f(x_k)|\geq k$。
我们取出点列 $\{x_k\}$，s.t. $|f(x_k)|\geq k$。直观上看
\[
\lim_{k\to+\infty}|f(x_k)|=+\infty .
\]

反过来看，如果存在点列 $\{x_k\}\subset X$ 使得
\[
\lim_{k\to+\infty}|f(x_k)|=+\infty,
\]
那么由极限定义
$\forall M>0,\ \exists N,\ \forall k>N,\ |f(x_k)|\geq M$，
可知 $f(x)$ 在 $X$ 上无界。

e.g. $f(x)=\dfrac1x$。取 $x_k=\dfrac1k$，$f(x_k)=k$，有
\[
\lim_{k\to+\infty}f(x_k)=+\infty,
\]
可知 $f$ 无界。

\section{函数}

\subsubsection*{复合函数}

\textbf{Def:}\quad $f:X\to\mathbb R$，$g:Y\to\mathbb R$。
如果 $f(X)\subseteq Y$，则 $g\circ f$ 有意义，定义为
\[
\forall x\in X,\quad (g\circ f)(x)=g\bigl(f(x)\bigr).
\]

如果 $f(X)\nsubseteq Y$，$g\circ f$ 也可以定义，
但需要把定义域调整为
\[
\{x\in X\mid f(x)\in Y\},\qquad\text{即}\quad D\triangleq f^{-1}(Y).
\]

\subsubsection*{反函数}

有函数 $f:X\to Y$。如果 $f$ 是单射，则可以定义逆映射
$f^{-1}:f(X)\to X$。

\textbf{Def:}\quad 当 $f:X\to\mathbb R$ 是函数且为单射，
$f^{-1}$ 称为\textbf{反函数}（Inverse function/map）。

e.g. $f:X\to\mathbb R$ 为严格单增函数，则 $f$ 存在反函数。

\textbf{pf:}\quad $\forall x_1\neq x_2$，不妨设 $x_1<x_2$，
则有 $f(x_1)<f(x_2)$，满足单射，存在反函数。

反过来，存在反函数不一定严格单调：
可以在几个区间上面严格单调，而在 $\mathbb R$ 上没有单调性。

\subsubsection*{初等函数}

指数函数、对数函数、幂函数、三角函数、多项式函数，
与它们通过加减乘除所生成的复合函数。

\section{极限}

\subsubsection*{极限的定义}

\begin{itemize}
\item \textbf{描述}\quad
\[
a=b\iff\forall\epsilon>0,\ |a-b|<\epsilon .
\]
\textbf{pf:}\quad 反设 $a\neq b$，不妨令 $a>b$，
取 $\epsilon_0=\dfrac{a-b}{2}$，则有
\[
|a-b|=a-b>\frac{a-b}{2}=\epsilon_0 .
\]

\item \textbf{描述趋势}\quad
现有数列 $\{a_n\}$，当 $n\to+\infty$ 时 $\{a_n\}$ 满足某性质 $P$，
即
\[
\exists N_P,\ \forall n>N_P\ \text{有}\ a_n\ \text{满足}\ P .
\]

\item \textbf{定义}\quad
$\forall\epsilon>0,\ \exists N_\epsilon,\ \forall n>N_\epsilon$ 有
$|a_n|<\epsilon$，则称 $a_n$ 的极限为 $0$，记为
\[
\lim_{n\to+\infty}a_n=0 .
\]

\item \textbf{从几何看极限}\quad 极限不存在。

对于数列 $\{S_n\}$，
\[
\{S_n\}=
\begin{cases}
n\ \text{是奇数}, & S_n=1,\\
n\ \text{是偶数}, & S_n=0 .
\end{cases}
\]
证明极限 $\lim\limits_{n\to+\infty}S_n$ 不存在。

分析：如果有极限，只可能是 $0$ 或 $1$，
我们证明两个都不是极限，说明极限不存在。下以不等于 $1$ 为例。

取 $\epsilon=\dfrac12$，对于任意的 $N$，总有比 $N$ 大的奇数，
使得 $S_n=1>\dfrac12$。

\item \textbf{趋于 $A$ 的极限}\quad
$\lim\limits_{n\to+\infty}a_n=A$，即
\[
\forall\epsilon>0,\ \exists N_\epsilon>0,\ \forall n>N_\epsilon,\
|a_n-A|<\epsilon .
\]

\item \textbf{无穷小量}\quad
\textbf{Def:}\quad $\lim\limits_{n\to+\infty}a_n=0$，
则称 $\{a_n\}$ 为无穷小量。

无穷小量不是指一个绝对值要多小有多小的常数，而是指的一个过程。
\end{itemize}

\end{document}

\end{document}
